% META: {"id": "TNSI-STRUCT-EX-002", "chapitre": "TNSI-STRUCTURES-DONNEES", "type_objet": "exercice", "capacites": ["T-STRUCT-02A", "T-STRUCT-02B"], "capacites_codes": ["C02A", "C02B"], "methodes": [], "parcours": 1, "competences": ["programmer"], "duree_min": 15, "mode_creation": "ex_nihilo", "sources_inspiration": [], "fichier_tex": "chapitres/TNSI-STRUCTURES-DONNEES/exercices/TNSI-STRUCT-EX-002.tex", "corrige_tex": "chapitres/TNSI-STRUCTURES-DONNEES/corriges/TNSI-STRUCT-CO-002.tex", "status": "generated"}
% BEGIN-VERIFY
% class Compte:
%     def __init__(self, titulaire, solde=0):
%         self.titulaire = titulaire
%         self.solde = solde
%     def deposer(self, montant):
%         self.solde = self.solde + montant
%     def retirer(self, montant):
%         if montant > self.solde:
%             raise ValueError('solde insuffisant')
%         self.solde = self.solde - montant
% c = Compte('Alice', 100)
% c.deposer(50)
% assert c.solde == 150
% c.retirer(30)
% assert c.solde == 120
% END-VERIFY
\begin{exercice}{TNSI-STRUCT-EX-002}{1}{15}
Ecrire une classe \lstinline{Compte} representant un compte bancaire simplifie, avec :
\begin{itemize}
  \item un constructeur prenant un \lstinline{titulaire} et un \lstinline{solde} initial (par defaut $0$) ;
  \item une methode \lstinline{deposer(montant)} qui augmente le solde ;
  \item une methode \lstinline{retirer(montant)} qui diminue le solde, et leve une erreur (\lstinline{ValueError}) si le montant demande depasse le solde disponible.
\end{itemize}
Tester votre classe avec un compte initial de $100$, un depot de $50$, puis un retrait de $30$.
\end{exercice}
